% !TeX root = ../main.tex

\chapter{Discussion}\label{chapter:discussion}

\section{Answers to the Research Questions}

\paragraph{RQ1: remaining trajectory accuracy.}
The accuracy difference from the closest higher-information reference is
reported directly. For an error measure $e$, let
\begin{equation}
  \Delta e_{\mathrm{abs}} = e_{\mathrm{PVO}}-e_{\mathrm{MV\text{-}first}},
  \qquad
  \Delta e_{\mathrm{rel}} =
  100\frac{e_{\mathrm{PVO}}-e_{\mathrm{MV\text{-}first}}}
  {e_{\mathrm{MV\text{-}first}}}.
  \label{eq:pvo-reference-error-difference}
\end{equation}
Positive values mean that PVO has greater error. For the easy group, the
differences between group macro means are \SI{0.334}{\metre} and 201.3\% for ATE, and
\SI{0.0101}{\metre} and 49.6\% for RTE-6. The small RTE-6 difference shows that
local motion can remain close to the MV-first reference on easier motion even
when the accumulated trajectory differs more strongly. Across the benchmark,
the gap grows and is dominated by several difficult sequences; EV203 is the
exception on which PVO has lower ATE and RTE-6 than MV-first. PVO nevertheless
improves on the stationary baseline in both metrics on nine of eleven
sequences. The relative differences must be read together with the absolute
ones because division by a small reference error can produce a large
percentage from a modest metric displacement.

\begin{table}[htbp]
  \centering
  \caption[Trajectory-error differences between PVO and VOCA MV-first]{Absolute and percentage differences between the group macro means
  of PVO and VOCA MV-first. Positive values indicate greater PVO error.
  Difficulty groups follow the EuRoC sequence labels.}
  \label{tab:pvo-mv-first-error-differences}
  \begin{tabular}{@{}lrrrr@{}}
    \toprule
    Group & $\Delta$ATE [m] & $\Delta$ATE [\%] & $\Delta$RTE-6 [m] & $\Delta$RTE-6 [\%] \\
    \midrule
    Easy & 0.334 & 201.3 & 0.0101 & 49.6 \\
    Medium & 1.091 & 796.6 & 0.0811 & 507.3 \\
    Difficult & 4.341 & 165.3 & 0.2954 & 72.4 \\
    All sequences & 1.998 & 189.7 & 0.1332 & 83.2 \\
    \bottomrule
  \end{tabular}
\end{table}

\paragraph{RQ2: stereo-temporal persistence and support.}
Persistent identities can be assembled without pixel descriptors, but their
typical lifetime is short. Across the benchmark, PVO confirms
5,230,916 tracks with a median length of three frames and a 95th percentile of
five frames. Only 9.92\% reach five frames. Stereo support is likewise uneven:
Between 13.48\% and 19.88\% of queries return a match on EMH01, EMH02, and EV101, whereas
the rate is at most 1.99\% on EV201 to EV203. In the separate structure
diagnostic, the proportion of codec supports that contain a Shi-Tomasi corner
changes across valid stereo, reciprocal agreement, robust registration, and
suppression: 30.59\% for raw vectors, 37.45\% with valid endpoint depth, 34.38\%
after reciprocal agreement, 39.57\% after TEASER++, and 44.54\% in the branched
suppression-retained pool. Robust registration therefore concentrates its input
near two-directional image structure, but most individual tracks still offer
little multiview support.

\paragraph{RQ3: properties associated with failure.}
The episode analysis identifies several distinct events
(\autoref{sec:failure-episodes}). EMH04 develops sustained drift while the
view becomes dark and sparsely detailed, despite abundant stereo-valid
candidates. EMH03 has a brief burst during rapid motion with substantially
fewer retained connections than another, slightly faster interval with similar
corner support. EV103 suffers a catastrophic pose update, temporary stationary
outputs, and recovery of local motion without recovery of the accumulated
trajectory. EV202 supplies both low-motion jumps and a failure in a
corner-rich view. Visible structure, observation availability, and usable
temporal geometry are therefore different quantities.

The paired residual comparison further limits the explanations. EMH03 and
EMH04 retain their selected failures with absolute reprojection, while EV103's
catastrophic error is much smaller. Another EV202 jump is worse with absolute
reprojection. The recorded events establish when support is lost and how the
estimated motion changes; they do not identify one cause shared by every
sequence. Track lifetime and consecutive-depth summaries provide context but
cannot explain an individual event without its temporal evidence.

\section{Operating Envelope}

The clearest feasibility evidence is the recovered trajectory geometry on the
easy sequences, particularly the Machine Hall views in
\autoref{sec:trajectory-geometry}. The reported PVO results on EMH01 and EV101
also show low accumulated error. Improvement over a stationary prediction on
all four easy sequences provides a broader check that the estimates contain
motion information. PVO also improves on all three medium sequences, although
the quality varies substantially. The EV202 margin is small: \SI{2.087}{\metre} ATE and
\SI{0.2078}{\metre} RTE-6 versus stationary values of \SI{2.111}{\metre} and
\SI{0.2441}{\metre}. The difficult set is mixed rather than uniformly
unsuccessful, with improvements on EMH04 and EMH05 but not on EV103 or EV203.

This grouping should not be read as a calibrated scale of codec difficulty.
The EuRoC labels summarize the original sequence challenge, not the quality of
H.264 prediction blocks or compressed-domain stereo. EV201 is labelled easy yet
has substantially lower stereo yield than the other easy sequences. Conversely,
local drift can remain moderate while small errors accumulate into a poor
global trajectory. The evidence supports an operating envelope for the fixed
EuRoC cameras, calibrations, and encoder configuration; it does not establish
the same boundary for other cameras, frame rates, codecs, or encoder searches.

\section{Track Persistence and Geometric Support}

The three-observation median is only one observation longer than the provisional
period required for confirmation. Consequently, most accepted landmarks supply
little temporal baseline before disappearing. This limits the number of views
that can jointly constrain a point and reduces the opportunity to average block
quantization and stereo-depth errors. The small tail of longer tracks shows that
identity persistence is possible, with a maximum of 204 frames, but that tail is
too sparse to characterize the typical backend input.

Stereo availability is one bottleneck rather than a complete explanation. A
continuation needs geometrically acceptable support, and a new identity also
needs depth at both endpoints before it can enter the robust 3D registration.
The low query yield on EV201 to EV203 therefore removes many potential
associations early. Conditional consecutive-depth changes are small among
points for which both depths exist, but this statistic says nothing about the
larger set with missing depth. Track loss can additionally reduce the connected
landmark ratio and request keyframes more frequently. Standard window
optimization and marginalization preserve information from accepted factors,
but they cannot recreate missing multiview observations.

The offline structure analysis provides a complementary result. Requiring valid
stereo at both endpoints, reciprocal agreement, and TEASER++ consensus changes
the fraction of regions containing a Shi-Tomasi corner at each stage. This
indicates whether geometric selection concentrates observations near
two-directional image structure without explicitly detecting it. It does not
show that each codec support contains one uniquely localized physical point,
and the diagnostic corner coordinates are never used by the estimator.

A quantitative lifetime comparison with the VOCA MV-first and OF-first
configurations is deliberately omitted. The diagnostic build instrumented for
track lifetimes did not reproduce their frozen trajectories, so its counters
would not describe the same configurations as the ATE and RTE-6 tables. The
qualitative expectation that pixel-refined tracking supplies
more precisely localized and longer-lived observations remains plausible, but
is not used here as a measured result.

\section{Error Modes}

The observed failures differ both in duration and in their relation to support.
EMH04 continues estimating motion with many observations but accumulates a large
orientation error. EV103's full-sequence RTE-6 is dominated by one catastrophic
window, followed by a period without emitted observations. These require
different explanations: the first calls into question the information supplied
by the surviving geometry, whereas the second also exposes the estimator's
response to an abrupt shortage of connections. Neither can be adequately
described by the full-sequence mean track length.

The observation-model ablation adds the same qualification at the backend.
Rigid flow lowers ATE on eight sequences, RTE-6 on seven, and both medians, but
its EV103 failure makes both macro means worse than absolute reprojection
(\autoref{tab:residual-ablation}). Replacing eligible absolute factors with
displacement factors lowers the median errors in this experiment, while
absolute initialization and fallback observations remain part of the estimator.
The replacement does not by itself guard against an erroneous correspondence
set or a poorly conditioned trajectory.

Two observation-level failures are directly visible before the later gates are
applied. First, local codec support can disappear while the tracked image
structure remains visible and raw temporal vectors remain elsewhere
(\autoref{fig:mv-abrupt-stop-example}). The responsible encoder decision is
unresolved in this example. Second, a codec block and its vector can
remain available while a localized appearance inside the source block is no
longer contained by the predicted destination block
(\autoref{fig:mv-feature-departure-example}). The latter is possible because the
encoder optimizes prediction of an extended block rather than persistence of a
particular scene point. These cases explain why raw-vector availability alone
does not establish a valid feature track.

In the separate EV103 episode, frame-type inspection identifies I-frames at the
relevant endpoints in both encoding directions, explaining the absence of ordinary
temporal vectors for that transition. Subsequent
P-frames restore raw candidates before useful frontend support returns. This
distinguishes an identified input interruption from the later losses whose
individual causes are not recorded. Similarly, EMH04's dark views suggest
ambiguous block predictions, but the saved logs cannot separate temporal-vector
bias from incorrect depth or an unfavorable distribution of points.

TEASER++ success does not by itself resolve those uncertainties. Its role is to
find a rigid transformation consistent with a subset of putative
associations~\cite{yang2020teaser}; such consistency does not independently
establish correspondence to the true scene. The implementation reports
successful registration throughout EMH04 even while the trajectory diverges.
Additional per-frame inlier geometry and residual diagnostics would be needed
to assess conditioning or identify coherent incorrect associations. The
present evidence does not establish that moving objects, calibration error,
or a particular depth distribution caused these episodes.

The median filter targets isolated compression artefacts, reciprocal agreement
tests temporal reversibility, TEASER++ rejects inconsistent 3D pairs, and
age-prioritized suppression removes indistinguishable duplicate support. These
stages address different symptoms. Their non-monotonic ablation results suggest
that stricter rejection can also discard useful constraints. The case studies
show where these distinctions matter, while identifying questions that the
available diagnostics cannot answer.

\section{Meaning of Pixel-Free Operation}

The information restriction applies to the tracking frontend and pose estimator. They
consume H.264 motion-vector side data, prediction-block geometry, timestamps,
camera calibration, rectification transforms, and validity masks. They do not
load decoded intensities, build image pyramids, detect corners, evaluate patches,
compute descriptors, or apply a photometric residual. Shi-Tomasi and blur
measurements are computed only after evaluation for diagnosis.

The restriction does not claim that the entire data-production process is
image-free. Decoded stereo images are remapped during offline rectification and
are then supplied to the video encoder. This conventional preprocessing creates
the H.264 streams on which the experiment operates. The resulting statement is
therefore precise: after calibrated offline rectification and encoding,
the frontend and pose estimator use codec side information and geometry without
consulting image intensities. The extraction software used here decodes frames
and reconstructs images internally, but does not forward those pixels to the
tracking frontend.

\section{Placement of Analysis and Transmission}\label{sec:cta-compatibility}

The communication boundary can be placed at several stages of the codec-based
frontend. Sending the required coded streams follows CTA: the capturing device
encodes, and the receiver extracts motion vectors and estimates camera motion.
Conventional CTA follows decoding with image-based feature extraction or
tracking~\cite{hofer2023h264cta}; here the frontend uses codec metadata directly.

Alternatively, the capturing device could extract the vectors and perform some
of the correspondence filtering before transmission. It could send temporal
and stereo observations for the remaining geometric checks, or run the frontend
locally and send confirmed observations to a remote backend. Analysis would then
precede the compression and transmission of the geometric data, giving an
ATC-like allocation of work while retaining codec-derived measurements. These
options extend the processing--communication choice described in the CTA--ATC
literature~\cite{redondi2016cta} to different stages of this frontend.

The transmitted representation must preserve the information needed after the
chosen boundary. Early motion-vector observations require block support,
reference information, timestamps, and calibrated geometry. A later split at
the frontend output would require shared calibration, consistent track identities,
per-camera coordinates, measured displacements, and predecessor timestamps,
with frontend and backend state maintained on their respective devices.
Filtering reduces the set of retained candidates, but the resulting bitrate also
depends on their metadata and encoding. The experiments do not implement or
measure these transport alternatives.

The evaluated input consists of two forward temporal streams, two reversed
temporal streams, and twelve shifted stereo streams. Preparing those additional
encodings has a cost. Two-pass temporal encoding and the current reversal of
complete frame lists make preparation an offline procedure, despite the
single-reference, B-frame-free prediction structure within each stream.
A live deployment must establish how the observations are generated and
delivered within its latency and throughput budget. Efficient metadata
extraction must also account for the pixel reconstruction performed by the
current decoder.

Passthrough+ motivates dedicated codec hardware with platform-specific power
budgets~\cite{chaurasia2020passthrough}. The present software-encoded experiments
establish geometric feasibility. Net energy savings depend on all preparation,
communication, extraction, and estimation. Reusing motion from an already
required stream has a different incremental cost from generating extra stereo
encodings. Sending an unchanged stream also preserves its bitrate; filtering
and transmitting geometric observations is a separate design choice.

\section{Validity and Limitations}

The evaluation covers all eleven EuRoC sequences, but only one benchmark, one
stereo camera family, and the provided calibrations. The result is also tied to
FFmpeg 4.4.2, \texttt{libx264}, H.264 High profile, the fixed temporal encoding,
and the twelve shifted I-P-P stereo streams. Hardware encoders or other search,
partition, bitrate, GOP, and colour-format choices may expose different motion
vectors. Passthrough+'s hardware mosaic, densification, pose-assisted temporal
depth propagation, and rendering stages are not reproduced. In particular, the
original full proxy assumes headset poses from an independent SLAM system,
which would violate the information restriction here. The stereo component is
therefore an adaptation rather than a complete replication of that system.

The selected PVO parameters were chosen after a bounded twelve-candidate sweep
on the same eleven sequences used for reporting. The final tables exclude
abandoned searches, but the evaluation is therefore not an independently
held-out benchmark study. The failure episodes are also selected retrospectively
and do not estimate the prevalence of a failure mechanism. Their successful
comparisons and counterexamples constrain the interpretation, but changing
views, illumination, motion, and available geometry remain coupled. The paired
residual ablation controls the observation model; image inspection and
timestamped counters alone are not equivalent to a controlled intervention.

The static-scene model, two-camera depth construction, calibrated
synchronization, and available ground truth further delimit the conclusions.
The system has no explicit mechanism for rolling shutter, independently moving
objects, calibration drift, or lost synchronization. A validity mask excludes
undefined rectification regions, but no image-derived confidence is available
inside valid regions.

\section{Comparison with Pixel-Refined References}

The VOCA MV-first configuration~\cite{hilscher2026voca} uses codec motion to
initialize KLT, followed by an unseeded optical-flow fallback. The VOCA OF-first
configuration attempts unseeded KLT first and uses motion-vector initialization
when tracks are lost. Both therefore receive information that the proposed
estimator intentionally omits.
Their median ATE values, \SI{0.200}{\metre} and \SI{0.118}{\metre}, are well
below PVO's \SI{1.082}{\metre}; their median RTE-6 values are
\SI{0.0163}{\metre} and \SI{0.0155}{\metre}, compared with
\SI{0.0818}{\metre}. The absolute and percentage differences from MV-first are
summarized in \autoref{tab:pvo-mv-first-error-differences}. This accuracy gap
remains substantial even though the MV-first reference fails on EV203.

These references contextualize the accuracy available with pixel-based
localization and refinement. PVO's nine improvements over the stationary baseline
demonstrate useful motion information under the pixel-free restriction. The
accuracy gap also includes differences in correspondence construction, filtering,
and observation models, so it does not isolate the effect of pixel refinement.
